\section{System Architecture Overview}


The designed system architecture for Phase~1 consists of completely isolated stream processing pipeline including message buffering, multi-level processing stages, and visualization of live streaming data. If the downstream analytics nodes are experiencing too much load and die, it will not affect the message buffering process in any way since it will be buffered through offsets in Kafka. Refer to Figure~\ref{fig:arch_diagram} for the separation of responsibilities.

\begin{figure}[htbp]
    \centering
    \includegraphics[width=0.85\textwidth]{Figures/architecture_diagram.png}
    \caption{Figure depicts a data pipeline, where there is a well-defined distinction between different functional modules, including event generators, Kafka brokers, analytical workers, Redis, FastAPI servers, and the graphical interface.}
    \label{fig:arch_diagram}
\end{figure}


The overall design consists of five key layers:
\begin{enumerate}
    \item \textbf{Stream Producer} --- Creates artificial events that follow a Zipfian distribution and sends them to a Kafka topic.
    \item \textbf{Kafka Message Broker} --- Provides the persistent and scalable messaging backbone that separates producers from consumers.
    \item \textbf{Analytics Workers} --- Autonomous worker threads (written in Python or C++), which ingest events, maintain sketch structures locally, and report statistics.
    \item \textbf{Redis Data Store} --- Distributed memory-backed database used to aggregate statistics collected by all workers.
    \item \textbf{FastAPI + Web Dashboard} --- Offers REST API services and a dynamic web-based dashboard for visualizing metrics in real-time.
\end{enumerate}

\section{Stream Ingestion Pipeline}

\subsection{Zipfian Traffic Generation}
A classic challenge in stream processing involves handling skewed traffic naturally. In practice, traffic exhibits power-law characteristics where a few items dominate the total volume of traffic~\cite{zipf}. The Stream Producer creates events using a Zipfian distribution that can be customized based on the skewness value (default $s = 1.2$). The events generated by the Stream Producer are then published in JSON format to the \texttt{stream\_events} topic of Kafka:

\begin{lstlisting}[language=Python, caption=Example format of the event published to Kafka.]
{"key": "item_42", "ts": 1710000000.123}
\end{lstlisting}

\subsection{Integration with Apache Kafka}
Apache Kafka~\cite{kafka} acts as the reliable event broker. The important design decisions include the following:
\begin{itemize}
    \item Each analytics worker consumes messages from Kafka using a \textbf{distinct consumer group}, making sure that all workers receive \textbf{all} events individually. It follows the analytics fanout model, wherein each node has its own copy of the sketch.
    \item Kafka offset mechanism offers inherent support for backpressure and fault tolerance.
    \item The producer supports the ability to configure event throughput rates (default: 5000 events/second) and stream lifetime.
\end{itemize}

\subsection{Concept Drift}
The production data is non-stationary, implying that popular objects are constantly changing. Thus, we require a way to accurately simulate such behavior in the experiments. For this purpose, we have provided a flag \texttt{--shift-interval}. Enabling this makes it so that the Zipfian rank distribution is changed periodically: the object which appeared most frequently is moved in favor of any random object from the end of the list.

\section{Worker Design for Analytics}

The workers asynchronously read events from their respective Kafka consumer groups. Each worker contains local instances of all three sketches and updates its metrics in Redis after certain intervals.

\subsection{Python Worker}
For a normal Python worker, the \texttt{kafka-python-ng} package is used for Kafka processing, and \texttt{hashlib.md5} is used for non-cryptographic hashing (extracting a 64-bit index from the digest). Some key aspects about this worker are:
\begin{itemize}
    \item Updates to sketches are done synchronously whenever an event arrives.
    \item P99 latency is calculated by taking the last 500 update timestamps into consideration.
    \item CPU and memory information is retrieved using \texttt{psutil}.
    \item Metrics are written to Redis at an interval of every 2 seconds.
\end{itemize}

\subsection{High-Performance C++ Worker}
In order to test performance ceilings theoretically, a high-performance C++20 drop-in replacement worker was created. It uses the same Kafka connection method (\texttt{librdkafka}) and writes to the same Redis keys, making it easy for the Python API, dashboard, and controller to run as is. Some benefits of its high performance include:

\begin{itemize}
    \item \textbf{XXH64 Hashing:} With a xxHash (XXH64) non-cryptographic hashing algorithm, there is approximately $10\times$ higher throughput than SHA-256, important due to the need for hash calculations for each event.
    \item \textbf{Bit Counting in Hardware:} \texttt{std::countl\_zero()} compiles down to a single hardware instruction (\texttt{LZCNT} or \texttt{BSR}) in HyperLogLog's leading-zero calculation, while in Python the equivalent while loop is used.
   \item \textbf{Parallel Reads:} The use of \texttt{std::shared\_mutex} enables parallel reading of the sketch without serializing execution due to the Global Interpreter Lock (GIL) in Python.
\end{itemize}

\section{Federated State Store}
Redis~\cite{redis} acts as the common ``blackboard'' among all system components. Table~\ref{tab:redis_keys} summarizes the key schema.

\begin{table}[htbp]
\centering
\begin{tabular}{l l l l}
\toprule
\textbf{Key Pattern} & \textbf{Contents} & \textbf{Written By} & \textbf{Read By} \\
\midrule
\texttt{worker:\{id\}:metrics}    & P99, mem\%, CMS error, HLL est. & Worker     & API, Dashboard \\
\texttt{worker:\{id\}:latency\_history} & Last 50 P99 samples       & Worker     & Dashboard \\
\texttt{worker:\{id\}:cmd}        & Controller commands             & Controller & Worker \\
\texttt{controller:log}           & Decision log                    & Controller & Dashboard \\
\texttt{router:query\_log}        & Routing decisions                & Router     & Dashboard \\
\bottomrule
\end{tabular}
\caption{Redis key schema used for inter-component communication.}
\label{tab:redis_keys}
\end{table}

Keys expire after 30 seconds TTL (time to live). In case a worker goes down or gets disconnected, keys expire and the worker gets kicked off the dashboard and scoring engine of the router.

\section{Analytics API}

Backend is developed using FastAPI~\cite{fastapi}, an advanced asynchronous framework written in Python. The API server handles both the REST endpoints as well as static dashboard pages. Key endpoints are provided in Table~\ref{tab:api_endpoints}.

\begin{table}[htbp]
\centering
\begin{tabular}{l l l}
\toprule
\textbf{Method} & \textbf{Endpoint} & \textbf{Description} \\
\midrule
GET  & \texttt{/query/frequency?key=X}   & Live routed CMS frequency estimate for key $X$ \\
GET  & \texttt{/query/cardinality}        & Routed HLL distinct count estimate \\
GET  & \texttt{/query/heavy-hitters?n=10} & Routed Top-$N$ Misra--Gries heavy hitters \\
GET  & \texttt{/query/distributed/frequency?key=X} & Scatter-gather Monoid-merged CMS frequency across all partitions \\
GET  & \texttt{/query/distributed/cardinality}      & Scatter-gather Monoid-merged HLL cardinality across all partitions \\
GET  & \texttt{/query/distributed/heavy-hitters?n=N}& Scatter-gather Space-Saving merged heavy hitters \\
GET  & \texttt{/metrics/cluster}          & Live health and telemetry for all active workers \\
GET  & \texttt{/metrics/partitions}       & Partition topology, primaries, and hedged replicas \\
GET  & \texttt{/metrics/controller-stats} & Live RL / Bandit step counts, $\epsilon$, and rewards \\
POST & \texttt{/control/mode}             & Switch controller mode (\texttt{rl}, \texttt{linucb}, \texttt{rule}) \\
POST & \texttt{/control/set-params}       & Force a precision template on a worker \\
GET  & \texttt{/health}                   & Liveness probe \\
\bottomrule
\end{tabular}
\caption{REST API endpoints implemented across Phase~1 and Phase~2.}
\label{tab:api_endpoints}
\end{table}

\section{Partitioned Ingestion Engine and Hedged Monoid Merge Architecture}
\label{sec:hedged_monoid_arch}

While the Phase~1 architecture relied on an analytics fan-out model (where every worker ingests the full stream), horizontally scaling write throughput to millions of events per second requires stream partitioning. In Phase~2, we developed a distributed \textbf{Partitioned Ingestion Engine with Hedged Monoid Merges}, enabling horizontal scaling while mitigating the straggler problem in distributed scatter-gather queries.

\subsection{Partition Topology and Replica Pairing}
The stream ingestion pipeline partitions events across $P$ disjoint Kafka partitions using deterministic key hashing:
\begin{equation}
    p = \text{CRC32}(k) \pmod P
\end{equation}
Each partition $p \in \{0, \dots, P-1\}$ is processed by a primary worker node $W_{p, \text{primary}}$ and backed by a hot standby replica $W_{p, \text{replica}}$. The primary worker actively ingests partition $p$ and answers incoming queries; the replica continuously mirrors partition $p$ in parallel. The Load-Aware Router discovers all active partitions via Redis heartbeats and dynamically designates primary and replica nodes for each partition based on the utility scoring function (Eq.~\ref{eq:node_score}).

\subsection{Monoid Merge Algebra for Distributed Sketches}
A fundamental property of probabilistic streaming data structures is that they form algebraic \textbf{Monoids} $(\mathcal{S}, \oplus, \mathbf{0})$, possessing an associative binary operation $\oplus$ and an identity element $\mathbf{0}$ such that:
\begin{align}
    (A \oplus B) \oplus C &= A \oplus (B \oplus C) \quad &&\text{(Associativity)} \\
    A \oplus \mathbf{0} &= \mathbf{0} \oplus A = A \quad &&\text{(Identity)}
\end{align}
Because raw stream events are partitioned disjointly across workers, a global query does not require pulling millions of raw events across the network. Instead, the scatter-gather engine retrieves tiny partial sketch summaries ($\sim 50\text{--}100\text{ KB}$) from the $P$ partitions and merges them into an exact global summary in $<500\ \mu\text{s}$:

\begin{enumerate}
    \item \textbf{Count-Min Sketch (Element-Wise Addition):}
    Let $\mathbf{C}_p \in \mathbb{N}^{d \times w}$ be the count matrix on partition $p$. The global merged sketch is the element-wise matrix sum:
    \begin{equation}
        \mathbf{C}_{\text{global}}[i, j] = \sum_{p=0}^{P-1} \mathbf{C}_p[i, j] \quad \forall i \in [d], j \in [w]
    \end{equation}
    Because addition in $\mathbb{N}$ is associative and commutative, Count-Min Sketch forms an Abelian Monoid with the all-zero matrix as identity. The point query estimate satisfies $\hat{f}(x) \le f_{\text{global}}(x) + \frac{e}{w} \sum_p N_p$.

    \item \textbf{HyperLogLog (Coordinate-Wise Register Maximum):}
    Let $\mathbf{R}_p \in \mathbb{N}^m$ be the register array of partition $p$ (where $m = 2^p$). The global register array is computed via coordinate-wise maximum:
    \begin{equation}
        \mathbf{R}_{\text{global}}[i] = \max_{p \in \{0, \dots, P-1\}} \mathbf{R}_p[i] \quad \forall i \in \{0, \dots, m-1\}
    \end{equation}
    Because $\max(a, b)$ is associative, commutative, and idempotent ($\max(a, a) = a$), HyperLogLog forms a bounded join-semilattice. Even if events are duplicated across partitions, the coordinate-wise maximum yields the exact deduplicated union cardinality across the cluster.

    \item \textbf{Misra--Gries (Space-Saving Monoid Merge):}
    To merge heavy-hitter summaries across partitions without exceeding memory budgets, we implement the Space-Saving Monoid merge algorithm (Agarwal et al., ACM TODS 2013). Candidate frequencies from all partitions are summed: $f_{\text{merged}}(x) = \sum_p f_p(x)$. If the number of unique candidate items exceeds the target capacity $k = \min_p k_p$, the $(k)^{\text{th}}$ largest count $C_{\min}$ is subtracted from all entries, retaining only candidates with positive remainder. The merged summary preserves bounded additive error:
    \begin{equation}
        \text{Error Bound} \le \frac{\sum_p N_p}{\min_p k_p}
    \end{equation}
\end{enumerate}

\subsection{Straggler Mitigation via Asynchronous Hedged Requests}
In any distributed scatter-gather query over $P$ partitions, the total query latency is strictly dictated by the slowest responder:
\begin{equation}
    T_{\text{query}} = \max_{p \in \{0, \dots, P-1\}} T_p
\end{equation}
In a cluster where individual nodes experience transient latency hiccups (such as Java/Python garbage collection pauses, operating system thread scheduling jitter, or background disk I/O) with even a small probability $p_{\text{strag}} = 5\%$, the probability of a scatter-gather query encountering at least one straggler escalates dramatically:
\begin{equation}
    P(\text{Query Stalls}) = 1 - (1 - p_{\text{strag}})^P
\end{equation}
For $P = 4$ partitions, $18.55\%$ of all queries stall. For $P = 16$ partitions, $56.0\%$ of all queries stall, completely destroying tail latency ($P95$ and $P99$).

To solve this, our \texttt{HedgedScatterGatherEngine} integrates \textbf{Hedged Requests} (inspired by Jeffrey Dean and Luiz André Barroso, \textit{The Tail at Scale}, CACM 2013):
\begin{enumerate}
    \item \textbf{Primary Dispatch:} The engine asynchronously sends the sub-query to the primary worker $W_{p, \text{primary}}$ for each partition $p$.
    \item \textbf{Hedge Timeout Check:} The engine sets a timer equal to $T_{\text{hedge}}$ (chosen as the empirical $P95$ normal response latency, typically $4.5\text{--}6.0\text{ ms}$).
    \item \textbf{Duplicate Hedged Dispatch:} If $W_{p, \text{primary}}$ does not respond within $T_{\text{hedge}}$, the engine immediately dispatches a duplicate hedged request to the replica worker $W_{p, \text{replica}}$.
    \item \textbf{First-Completion Acceptance and Cancellation:} Whichever node responds first provides the result for partition $p$; the outstanding task on the straggler node is cancelled.
\end{enumerate}
The effective latency for partition $p$ becomes:
\begin{equation}
    T_p^{\text{hedged}} = \min\left(T_{p, \text{primary}},\; T_{\text{hedge}} + T_{p, \text{replica}}\right)
\end{equation}
Because the probability of both the primary and the replica simultaneously suffering a straggler stall is $p_{\text{strag}}^2 = (0.05)^2 = 0.25\%$, tail latency collapses back down to normal operational bounds while introducing only $\sim 5\%$ additional network request overhead.
